\section{Lựa chọn và huấn luyện mô hình}

\subsection{Chuẩn bị dữ liệu cho mô hình}
\begin{frame}
    \textbf{Phân chia tập dữ liệu chống rò rỉ}\\
    \vspace{0.1cm}
    Tập dữ liệu sạch gồm 7.149 ảnh được phân chia thành ba tập theo tỷ lệ chuẩn $70:15:15$ bằng kỹ thuật Stratified Sampling.

    \vspace{0.2cm}
    \begin{table}[h]
        \centering
        \caption{Phân bố số lượng ảnh theo tập con, miền và lớp bệnh}
        \label{tab:slide_split_dist}
        \footnotesize
        \setlength{\tabcolsep}{10pt}
        \begin{tabular}{lrrrr}
            \toprule
            Miền dữ liệu                   & Train (70\%)   & Val (15\%)     & Test (15\%)    & Tổng cộng      \\
            \midrule
            Lúa                            & 2.346          & 502            & 503            & 3.351          \\
            Cà phê                         & 2.658          & 570            & 570            & 3.798          \\
            \midrule
            \textbf{Tổng toàn cục (8 lớp)} & \textbf{5.004} & \textbf{1.072} & \textbf{1.073} & \textbf{7.149} \\
            \bottomrule
        \end{tabular}
    \end{table}
\end{frame}

\begin{frame}
    \textbf{Tiền xử lý và tăng cường dữ liệu}\\
    \footnotesize
    \justifying
    Để tăng khả năng tổng quát hóa và chống overfitting, nhóm áp dụng các kỹ thuật chuẩn hóa và tăng cường dữ liệu trong quá trình huấn luyện:

    \begin{table}[h]
        \centering
        \caption{Các kỹ thuật tiền xử lý và tăng cường dữ liệu}
        \label{tab:slide_data_aug}
        \scriptsize
        \setlength{\tabcolsep}{6pt}
        \begin{tabular}{lp{9.8cm}}
            \toprule
            Phương pháp               & Chi tiết thiết lập và tác dụng                                                                                              \\
            \midrule
            Letterbox                 & Đưa ảnh về kích thước $640 \times 640$ hoặc $256 \times 256$ px giữ nguyên tỷ lệ khung hình thực, tránh méo dạng vùng bệnh. \\
            Cắt và đổi kích thước     & Cắt và thay đổi kích thước ngẫu nhiên với tỷ lệ diện tích từ $0,75$ đến $1,0$.                                              \\
            Lật hình học              & Lật ngang với xác suất $p=0,5$, lật dọc với xác suất $p=0,2$.                                                               \\
            Biến đổi affine           & Xoay ngẫu nhiên trong khoảng góc $\pm 30^\circ$.                                                                            \\
            Biến thiên quang học      & Điều chỉnh ngẫu nhiên độ sáng và màu sắc ảnh.                                                                               \\
            Trộn mẫu (\textit{MixUp}) & Mixup với hệ số $\alpha=0,4$.                                                                                               \\
            \bottomrule
        \end{tabular}
    \end{table}
\end{frame}

\subsection{Lựa chọn và kiến trúc mô hình}
\begin{frame}
    \footnotesize
    \justifying
    Nhóm khảo sát, huấn luyện và so sánh năm kiến trúc tiên tiến thuộc hai trường phái Convolutional Neural Network (CNN) và Vision Transformer (ViT) nhằm tối ưu điểm cân bằng giữa độ chính xác và tốc độ suy luận:

    \vspace{0.1cm}
    \begin{table}[h]
        \centering
        \caption{Tổng quan năm kiến trúc cho bài toán instance segmentation}
        \label{tab:slide_model_overview}
        \scriptsize
        \setlength{\tabcolsep}{4pt}
        \begin{tabular}{llp{9.4cm}}
            \toprule
            Mô hình                                                 & Họ  & Đặc điểm nổi bật và lý do lựa chọn                                                                             \\
            \midrule
            Mask R-CNN \cite{he2017mask}                            & CNN & Kiến trúc hai giai đoạn kinh điển, mang lại độ chính xác cao, được chọn làm baseline đối chiếu hiệu suất.      \\
            YOLO26-seg \cite{jocher2023ultralytics,jocher2024yolo} & CNN & Kiến trúc một giai đoạn với tốc độ thời gian thực xuất sắc, cỡ medium tối ưu điểm cân bằng chính xác - tốc độ. \\
            RF-DETR \cite{carion2020end,lv2023detrs}                & ViT & Thiết kế end-to-end loại bỏ NMS, DINOv2 feature extraction, real-time Transformer.                             \\
            MobileSAM \cite{zhang2023faster,kirillov2023segment}    & ViT & Phiên bản thu gọn siêu nhẹ của SAM, khả năng zero-shot với câu lệnh đầu vào, tối ưu cho edge devices.          \\
            Mask2Former \cite{cheng2022masked}                      & ViT & Kiến trúc phân vùng hợp nhất dựa trên mask classification, đạt hiệu năng đa tác vụ vượt trội.                  \\
            \bottomrule
        \end{tabular}
    \end{table}
\end{frame}

\begin{frame}{Kiến trúc mô hình YOLO26-seg}
    \begin{columns}[T]
        \begin{column}{0.45\textwidth}
            \justifying
            \centerline{\textbf{Cấu trúc phân tầng của mô hình}}
            \vspace{0.2cm}
            
            YOLO26-seg sử dụng mạng tích chập sâu với ba khối chức năng chính:
            \begin{itemize}
                \setlength{\itemsep}{1pt}
                \item[-] \textbf{Backbone}: Trích xuất đa quy mô đặc trưng ngữ cảnh từ hình ảnh lá bằng các khối C2f và SPPF.
                \item[-] \textbf{Neck}: Kết hợp luồng thông tin PANet để tổng hợp các đặc trưng không gian và ngữ nghĩa cấp cao.
                \item[-] \textbf{Head}: Nhánh dự đoán song song (Decoupled Head) tách biệt cho hộp bao và mặt nạ phân vùng.
            \end{itemize}
        \end{column}
        \begin{column}{0.51\textwidth}
            \centering
            \resizebox{\textwidth}{!}{
            \begin{tikzpicture}[
                node distance=0.8cm and 0.4cm,
                block/.style={rectangle, draw, fill=Logo1!10, text width=2.8cm, text centered, rounded corners, minimum height=1.0cm, font=\small},
                branch/.style={rectangle, draw, fill=Logo2!10, text width=3.2cm, text centered, rounded corners, minimum height=1.0cm, font=\small},
                outblock/.style={rectangle, draw, fill=Logo1!20, text width=2.8cm, text centered, rounded corners, minimum height=1.0cm, font=\small},
                line/.style={draw, -{Stealth[scale=1.2]}, thick}
            ]
                % Nodes
                \node [block] (input) {Ảnh đầu vào \\ $640 \times 640 \times 3$};
                \node [block] (backbone) [below=of input] {Backbone \\ (C2f, SPPF)};
                \node [block] (neck) [below=of backbone] {Neck (PANet) \\ (FPN + PAN)};
                
                % Branches
                \node [branch] (boxhead) [below left=0.6cm and 0.1cm of neck] {Nhánh dự đoán hộp bao \\ (Box \& Class)};
                \node [branch] (maskhead) [below right=0.6cm and 0.1cm of neck] {Nhánh dự đoán mặt nạ \\ (Proto \& Coefficients)};
                
                \node [outblock] (output) [below=1.8cm of neck] {Kết quả phân vùng \\ (Segmentation)};
            
                % Lines
                \path [line] (input) -- (backbone);
                \path [line] (backbone) -- (neck);
                \path [line] (neck.south) -- ++(0,-0.3) -| (boxhead.north);
                \path [line] (neck.south) -- ++(0,-0.3) -| (maskhead.north);
                \path [line] (boxhead.south) -- ++(0,-0.3) -| (output.west);
                \path [line] (maskhead.south) -- ++(0,-0.3) -| (output.east);
            \end{tikzpicture}
            }
            \captionof{figure}{Sơ đồ kiến trúc YOLO26-seg}
        \end{column}
    \end{columns}
\end{frame}

\begin{frame}{Hàm mất mát của YOLO26-seg}
    Bài toán phân vùng thực thể vết bệnh được tối ưu hóa thông qua hàm mất mát tổng hợp:
    \begin{equation}
        \mathcal{L}_{total} = \lambda_{box} \mathcal{L}_{box} + \lambda_{cls} \mathcal{L}_{cls} + \lambda_{dfl} \mathcal{L}_{dfl} + \lambda_{mask} \mathcal{L}_{mask}
    \end{equation}
    \vspace{0.2cm}
    Trong đó:
    \begin{itemize}
        \setlength{\itemsep}{3pt}
        \item[-] $\mathcal{L}_{box}$ (Complete IoU Loss): Đo lường sai lệch vị trí, tỷ lệ khung hình và kích thước hộp bao.
        \item[-] $\mathcal{L}_{cls}$ (Binary Cross Entropy): Phân loại nhãn lớp bệnh cho từng thực thể.
        \item[-] $\mathcal{L}_{dfl}$ (Distribution Focal Loss): Tinh chỉnh biên hộp bao trong các trường hợp ranh giới mờ nhòe.
        \item[-] $\mathcal{L}_{mask}$ (BCE + Dice Loss): Tối ưu hóa sự trùng khớp mặt nạ phân vùng ở cấp độ điểm ảnh.
    \end{itemize}
\end{frame}

\subsection{Cấu hình huấn luyện}
\begin{frame}
    \begin{columns}[T]
        \begin{column}{0.44\textwidth}
            \justifying
            \centerline{\textbf{Thiết lập cấu hình chung}}
            \vspace{0.2cm}
            \begin{table}[h]
                \centering
                \caption{Cấu hình huấn luyện chung}
                \label{tab:slide_common_train}
                \footnotesize
                \setlength{\tabcolsep}{4pt}
                \begin{tabular}{ll}
                    \toprule
                    Tham số              & Giá trị thiết lập \\
                    \midrule
                    Phần cứng GPU        & NVIDIA Tesla T4   \\
                    Tỷ lệ chia dữ liệu   & $70:15:15$        \\
                    Thuật toán tối ưu    & AdamW             \\
                    Effective batch size & 8                 \\
                    Số epoch tối đa      & 30                \\
                    Dừng sớm             & 5                 \\
                    Seed                 & 3407              \\
                    \bottomrule
                \end{tabular}
            \end{table}
        \end{column}
        \begin{column}{0.54\textwidth}
            \justifying
            \centerline{\textbf{Tham số riêng theo kiến trúc}}
            \vspace{0.2cm}
            \begin{table}[h]
                \centering
                \caption{Cấu hình riêng cho từng mô hình}
                \label{tab:slide_specific_train}
                \scriptsize
                \setlength{\tabcolsep}{3pt}
                \begin{tabular}{lrrp{2.6cm}}
                    \toprule
                    Mô hình     & $\text{lr}_0$     & Kích thước         & Ghi chú kỹ thuật    \\
                    \midrule
                    Mask R-CNN  & $1 \cdot 10^{-4}$ & $256 \times 256$   & Batch 2, grad acc 4 \\
                    YOLO26-seg  & $8 \cdot 10^{-4}$ & $640 \times 640$   & Batch 8             \\
                    RF-DETR     & $1 \cdot 10^{-4}$ & $640 \times 640$   & Batch 2, grad acc 4 \\
                    MobileSAM   & $1 \cdot 10^{-5}$ & $1024 \times 1024$ & Batch 2, grad acc 4 \\
                    Mask2Former & $5 \cdot 10^{-5}$ & $384 \times 384$   & Batch 4, grad acc 2 \\
                    \bottomrule
                \end{tabular}
            \end{table}
        \end{column}
    \end{columns}
    \vspace{0.1cm}
    \scriptsize
    \textit{*Ghi chú: Batch: batch size, grad acc: tích lũy đạo hàm.}
\end{frame}


\subsection{Tinh chỉnh siêu tham số}
\begin{frame}
    \begin{columns}[T]
        \begin{column}{0.48\textwidth}
            \justifying
            \centerline{\textbf{Tối ưu tự động bằng Ray Tune}}
            \vspace{0.2cm}
            \footnotesize
            \begin{itemize}
                \setlength{\itemsep}{3pt}
                \item Tinh chỉnh siêu tham số cho mô hình tốt nhất YOLO26-seg bằng Ray Tune.
                \item Áp dụng thuật toán lập lịch bất đồng bộ ASHA.
                \item Đánh giá song song nhiều cấu hình thử nghiệm, dừng sớm các cấu hình kém hiệu quả ngay từ những epoch đầu giúp tối ưu hóa tài nguyên tính toán.
                \item Thực nghiệm 20 thử nghiệm độc lập trên không gian sáu siêu tham số cốt lõi.
            \end{itemize}
        \end{column}
        \begin{column}{0.48\textwidth}
            \centering
            \vspace{0.15cm}
            \begin{table}[h]
                \centering
                \caption{Không gian tìm kiếm siêu tham số}
                \label{tab:slide_raytune_space}
                \scriptsize
                \setlength{\tabcolsep}{3pt}
                \begin{tabular}{llr}
                    \toprule
                    Siêu tham số                           & Phân phối   & Khoảng/Tập giá trị                   \\
                    \midrule
                    Learning rate ($\text{lr}_0$)          & Log-uniform & $[1 \cdot 10^{-5}, 1 \cdot 10^{-2}]$ \\
                    Momentum                               & Uniform     & $[0,85, 0,95]$                       \\
                    Weight decay                           & Log-uniform & $[1 \cdot 10^{-5}, 1 \cdot 10^{-3}]$ \\
                    Warmup epochs                          & Choice      & $\{1, 2, 3\}$                        \\
                    Hệ số khung bao $\lambda_{\text{box}}$ & Uniform     & $[5,0, 10,0]$                        \\
                    Hệ số phân loại $\lambda_{\text{cls}}$ & Uniform     & $[0,3, 0,8]$                         \\
                    \bottomrule
                \end{tabular}
            \end{table}
        \end{column}
    \end{columns}
\end{frame}
